\section{Hipótese}
Parte-se da hipótese de que alunos submetidos à metodologia expositiva em sala de aula virtual tendem a apresentar dispersão de foco e baixo engajamento, em razão das múltiplas distrações presentes em seu ambiente doméstico e da limitação do docente em monitorar e intervir sobre comportamentos que ocorrem de forma velada por trás da tela.

A formulação de um problema solucionável é fundamental para o início de uma pesquisa científica, sendo a etapa seguinte a proposição de uma solução, denominada hipótese <citar-gil>.

Para este trabalho, propõe-se uma abordagem de associação entre variáveis, na qual as métricas coletadas serão associadas a variáveis representativas de atenção e engajamento. Serão observados, respectivamente, o tempo fora da conferência (aula) e o estado dos periféricos. <citar-gil> afirma que esse tipo de hipótese busca indicar a força ou o sentido da relação variável-valor, mas não determina causa, dependência ou influência.